\documentclass[aspectratio=169]{beamer}
\usetheme{Madrid}
\usecolortheme{default}
\usepackage[utf8]{inputenc}
\usepackage[spanish]{babel}
\usepackage{graphicx}
\usepackage{hyperref}

\title{Normas APA - Séptima Edición}
\subtitle{Introducción y Generalidades}
\author{Edison Achalma}
\institute{Universidad Nacional de San Cristóbal de Huamanga}
\date{\today}

\begin{document}

\frame{\titlepage}

\begin{frame}{Tabla de Contenido}
\tableofcontents
\end{frame}

\section{¿Qué son las Normas APA?}

\begin{frame}{¿Qué son las Normas APA?}
\begin{itemize}
    \item Conjunto de estándares creados por la \textbf{American Psychological Association}
    \item Utilizadas para la presentación de trabajos académicos y científicos
    \item Facilitan la comunicación clara y precisa en las ciencias sociales
    \item Actualmente en su \textbf{séptima edición} (2019)
\end{itemize}
\vspace{0.5cm}
% Imagen: Portada del manual APA 7ma edición
\begin{center}
\includegraphics[width=0.4\textwidth]{images/portada_apa7.png}
\end{center}
\end{frame}

\begin{frame}{Historia y Evolución}
\begin{columns}
\column{0.5\textwidth}
\begin{itemize}
    \item Primera edición: 1929
    \item Séptima edición: 2019
    \item Uso extendido en:
    \begin{itemize}
        \item Ciencias Sociales
        \item Humanidades
        \item Ciencias Administrativas
        \item Ciencias de la Salud
    \end{itemize}
\end{itemize}

\column{0.5\textwidth}
% Imagen: Línea de tiempo de las ediciones APA
\includegraphics[width=\textwidth]{images/historia_apa.png}
\end{columns}
\end{frame}

\section{Áreas de Aplicación}

\begin{frame}{Campos que Usan Normas APA}
% Imagen: Diagrama circular con las disciplinas que usan APA
\begin{center}
\includegraphics[width=0.8\textwidth]{images/disciplinas_apa.png}
\end{center}
\end{frame}

\begin{frame}{¿Por qué usar Normas APA?}
\begin{block}{Ventajas}
\begin{itemize}
    \item Estandarización de trabajos académicos
    \item Facilita la localización de fuentes
    \item Otorga crédito apropiado a los autores
    \item Mejora la credibilidad del trabajo
    \item Evita el plagio académico
\end{itemize}
\end{block}

\begin{alertblock}{Importante}
La correcta aplicación de las normas APA es fundamental en la integridad académica
\end{alertblock}
\end{frame}

\section{Elementos Principales}

\begin{frame}{Componentes de las Normas APA}
Las Normas APA abarcan cuatro aspectos principales:
\begin{enumerate}
    \item \textbf{Estructura} del trabajo
    \item \textbf{Formato} general del documento
    \item \textbf{Estilo} de redacción
    \item \textbf{Citación} y referencias
\end{enumerate}

\vspace{0.5cm}
% Imagen: Esquema de los 4 componentes principales
\begin{center}
\includegraphics[width=0.7\textwidth]{images/componentes_apa.png}
\end{center}
\end{frame}

\begin{frame}{Novedades de la 7ª Edición}
\begin{columns}
\column{0.5\textwidth}
\textbf{Principales cambios:}
\begin{itemize}
    \item Nuevas fuentes aceptadas
    \item Simplificación de formatos
    \item Eliminación de "Recuperado de"
    \item Uso de DOI simplificado
    \item Formato para redes sociales
\end{itemize}

\column{0.5\textwidth}
% Imagen: Comparación APA 6 vs APA 7
\includegraphics[width=\textwidth]{images/cambios_apa6_apa7.png}
\end{columns}
\end{frame}

\section{Recursos Oficiales}

\begin{frame}{Recursos para Consulta}
\begin{block}{Sitio Web Oficial}
\url{https://apastyle.apa.org}
\end{block}

\begin{block}{Manual Oficial}
Publication Manual of the American Psychological Association (7th ed.)
\end{block}

% Imagen: Captura del sitio web de APA Style
\begin{center}
\includegraphics[width=0.6\textwidth]{images/sitio_web_apa.png}
\end{center}
\end{frame}

\begin{frame}{Abreviaturas Comunes en APA}
% Imagen: Tabla de abreviaturas más utilizadas (cap., ed., p., pp., Vol., etc.)
\begin{center}
\includegraphics[width=0.9\textwidth]{images/abreviaturas_apa.png}
\end{center}
\end{frame}

\begin{frame}[standout]
\Huge ¿Preguntas?
\end{frame}

\begin{frame}{Próxima Clase}
\begin{center}
\Large \textbf{Estructura del Trabajo Académico}
\vspace{1cm}

Veremos:
\begin{itemize}
    \item Página de presentación
    \item Resumen y palabras clave
    \item Cuerpo del texto
    \item Referencias y anexos
\end{itemize}
\end{center}
\end{frame}

\end{document}
